{\footnotesize
\begin{description}[style=unboxed,leftmargin=0.6cm,font=\normalfont\bfseries]

\item[European Union] Europe; Advanced; overall 9.1/10; European Medicines Agency (EMA), National Competent Authorities, Notified Bodies; 600 devices; first AI instrument 2018. Scores: P10 C9 A9 T9 E10 PM9 L8. Privacy: GDPR (General Data Protection Regulation). AI rule: EU AI Act (2024). Challenges: Complex multi-layered regulatory landscape. Harmonization across 27 member states. Notified Body capacity constraints. Risk of over-regulation stifling innov... Notable: EU AI Act is the world's first comprehensive AI law. High-risk healthcare AI must comply by 2026-2027. EHDS aims to create a unified health data framework.
\item[Germany] Europe; Advanced; overall 8.7/10; BfArM (Federal Institute for Drugs and Medical Devices); 180 devices; first AI instrument 2018. Scores: P10 C9 A9 T8 E9 PM9 L7. Privacy: GDPR + BDSG (Federal Data Protection Act). AI rule: EU AI Act (as EU member). National AI Strategy (2018, updated 2020).. Challenges: Historically strict data protection culture can slow health data utilization. DiGA program growing pains. Interoperability of health IT systems. Notable: DiGA is the world's first regulatory pathway for prescribed digital therapies reimbursed by statutory health insurance. Health Data Use Act aims to unlock re...
\item[United States] North America; Advanced; overall 7.6/10; FDA (Food and Drug Administration); 950 devices; first AI instrument 2017. Scores: P8 C9 A9 T6 E6 PM8 L7. Privacy: HIPAA (Health Insurance Portability and Accountability Act). AI rule: FDA AI/ML-Based SaMD Action Plan (2021). Challenges: Fragmented state-level privacy laws. No comprehensive federal AI legislation. Balancing innovation speed with safety oversight for adaptive algorithms. Notable: FDA has authorized over 950 AI/ML-enabled medical devices as of 2025. Predetermined Change Control Plan allows manufacturers to pre-specify algorithm modific...
\item[United Kingdom] Europe; Advanced; overall 7.1/10; MHRA (Medicines and Healthcare products Regulatory Agency); 200 devices; first AI instrument 2018. Scores: P8 C8 A7 T7 E7 PM7 L6. Privacy: UK GDPR \& Data Protection Act 2018. AI rule: Pro-innovation AI Regulation Framework (2023). Challenges: Post-Brexit regulatory divergence from EU. Balancing light-touch innovation-friendly approach with patient safety. Transition from CE to UKCA marking. Notable: UK positioned as pro-innovation alternative to EU AI Act. AI Safety Institute is a world-first government body for AI safety research. NHS AI Lab supports AI...
\item[Canada] North America; Advanced; overall 7.1/10; Health Canada; 120 devices; first AI instrument 2019. Scores: P7 C8 A8 T7 E7 PM7 L6. Privacy: PIPEDA (Personal Information Protection and Electronic Documents Act), Provincial health privacy laws (e.g., Ontario PHIPA). AI rule: AIDA (Artificial Intelligence and Data Act, proposed in Bill C-27). Challenges: Federal-provincial jurisdiction complexity. AIDA legislation delayed. Keeping pace with rapid AI advancement in healthcare. Notable: Canada co-developed GMLP with FDA and MHRA. Pan-Canadian Health Data Strategy in development. Strong AI research ecosystem (MILA, Vector, Amii) informing pol...
\item[Japan] Asia; Advanced; overall 7.0/10; PMDA (Pharmaceuticals and Medical Devices Agency); 150 devices; first AI instrument 2019. Scores: P7 C8 A8 T6 E7 PM8 L5. Privacy: APPI (Act on Protection of Personal Information, amended 2022). AI rule: AI Guidelines for Business (2024), Social Principles of Human-Centric AI (2019). Challenges: Aging society drives demand but conservative regulatory culture. Language barriers in international harmonization. Slow adoption of digital health relative t... Notable: Japan's regulatory sandbox (IDATEN) accelerates innovative device approval. Society 5.0 vision integrates AI in healthcare transformation. PMDA actively deve...
\item[Singapore] Asia; Advanced; overall 6.9/10; HSA (Health Sciences Authority); 60 devices; first AI instrument 2019. Scores: P7 C7 A7 T8 E8 PM6 L5. Privacy: PDPA (Personal Data Protection Act, 2012, amended 2021). AI rule: Model AI Governance Framework (2019, updated 2020), AI Verify (2022). Challenges: Small domestic market limits local clinical evidence generation. Balancing business-friendly environment with robust oversight. Regional influence vs. global... Notable: AI Verify is the world's first AI governance testing framework. Singapore positioned as AI governance thought leader in ASEAN. Strong public-private partners...
\item[Switzerland] Europe; Advanced; overall 6.9/10; Swissmedic; 90 devices; first AI instrument 2023. Scores: P8 C8 A8 T6 E6 PM7 L5. Privacy: nFADP (New Federal Act on Data Protection, 2023). AI rule: No AI-specific law. Federal Council AI guidelines. Bilateral alignment with EU standards.. Challenges: Dependency on EU regulatory decisions without EU membership voting rights. Maintaining bilateral agreements post-EU AI Act. Small market size. Notable: Strong pharma/medtech hub (Novartis, Roche). nFADP brings Swiss privacy law to GDPR-equivalent. Active AI health research at ETH Zurich and EPFL.
\item[China] Asia; Advanced; overall 6.7/10; NMPA (National Medical Products Administration); 300 devices; first AI instrument 2021. Scores: P7 C7 A7 T7 E6 PM7 L6. Privacy: PIPL (Personal Information Protection Law, 2021). AI rule: Interim Measures for Generative AI (2023), Algorithm Recommendation Regulations (2022), Deep Synthesis Regulations (2023). Challenges: Rapid regulatory evolution creates compliance complexity. Data localization requirements limit international collaboration. Balancing state surveillance prio... Notable: China is among the first to regulate algorithmic recommendations and generative AI. NMPA rapidly approving AI diagnostic tools. National AI standardization e...
\item[Australia] Oceania; Moderate; overall 6.6/10; TGA (Therapeutic Goods Administration); 80 devices; first AI instrument 2019. Scores: P7 C7 A7 T6 E7 PM7 L5. Privacy: Privacy Act 1988 (amended), Australian Privacy Principles (APPs). AI rule: Australia's AI Ethics Principles (2019), Voluntary AI Safety Standard (2024). Challenges: Voluntary nature of AI ethics framework limits enforcement. Resource constraints in TGA for AI evaluation. Geographic barriers to clinical trial recruitment. Notable: Australia's AI Ethics Principles among the earliest globally. TGA actively developing SaMD regulatory pathways. National Digital Health Strategy driving AI i...
\item[South Korea] Asia; Advanced; overall 6.6/10; MFDS (Ministry of Food and Drug Safety); 200 devices; first AI instrument 2019. Scores: P7 C7 A8 T6 E6 PM7 L5. Privacy: PIPA (Personal Information Protection Act, amended 2023). AI rule: AI Framework Act (proposed), National AI Strategy (2019). Challenges: Balancing rapid technological adoption with regulatory oversight. Harmonizing domestic standards with international frameworks. SME compliance burden. Notable: South Korea has one of the highest AI medical device approval rates globally. MFDS Regulatory Science Center is a dedicated AI device evaluation unit. K-BioH...
\item[Israel] Middle East; Advanced; overall 6.3/10; MOH AMAR (Division of Medical Devices, Ministry of Health); 100 devices; first AI instrument 2022. Scores: P6 C8 A7 T6 E6 PM6 L5. Privacy: Privacy Protection Law (1981, amended), Health Information Exchange regulations. AI rule: AI Policy and Ethics Framework (2022), National AI Program. Challenges: Outdated base privacy legislation. Balancing startup-friendly ecosystem with regulatory rigor. Small domestic market necessitates international focus. Notable: Israel is a global health AI innovation hub. Unique national health data infrastructure enables world-class AI research. Highest health AI startups per capit...
\item[Brazil] South America; Developing; overall 5.3/10; ANVISA (Agência Nacional de Vigilância Sanitária); 40 devices; first AI instrument 2020. Scores: P7 C5 A5 T5 E5 PM5 L5. Privacy: LGPD (Lei Geral de Proteção de Dados, 2020). AI rule: AI Bill (PL 2338/2023, under legislative process). Challenges: AI regulatory framework still in legislative process. Limited ANVISA capacity for AI device evaluation. Digital health infrastructure gaps in underserved reg... Notable: LGPD is Latin America's most comprehensive data protection law. Active AI legislation process. Growing health AI startup ecosystem in São Paulo.
\item[UAE] Middle East; Emerging; overall 4.6/10; MOH (Ministry of Health and Prevention), HAAD, DHA; 35 devices; first AI instrument 2017. Scores: P5 C5 A5 T5 E5 PM4 L3. Privacy: Federal Decree-Law No. 45/2021 (Personal Data Protection). AI rule: National AI Strategy 2031, Minister of State for AI appointed (2017). Challenges: Multi-authority regulatory landscape (federal and emirate-level). Rapid strategy ambitions vs. regulatory maturity. Building domestic healthcare data infrast... Notable: First country to appoint a Minister of State for AI (2017). Ambitious National AI Strategy 2031. Hub for health AI startups and international collaborations.
\item[Saudi Arabia] Middle East; Emerging; overall 4.4/10; SFDA (Saudi Food and Drug Authority); 25 devices; first AI instrument 2023. Scores: P5 C4 A5 T5 E5 PM4 L3. Privacy: PDPL (Personal Data Protection Law, 2023). AI rule: SDAIA AI Ethics Principles (2023), National AI Strategy. Challenges: Regulatory framework still maturing. Building domestic regulatory expertise. Dependency on international recognition pathways. Limited public health data inf... Notable: Massive AI investment under Vision 2030. SDAIA established as dedicated AI authority. NEOM project integrating AI healthcare. Rapid digital health transforma...
\item[India] Asia; Developing; overall 4.3/10; CDSCO (Central Drugs Standard Control Organisation); 30 devices; first AI instrument 2023. Scores: P5 C4 A4 T4 E5 PM4 L4. Privacy: Digital Personal Data Protection Act (DPDP, 2023). AI rule: No AI-specific law. NITI Aayog AI Strategy and Responsible AI guidelines (2021).. Challenges: Regulatory framework still maturing for AI devices. Limited CDSCO capacity for AI/ML evaluation. Infrastructure gaps in digital health. Balancing innovation ... Notable: DPDP Act 2023 is India's first comprehensive data protection law. Ayushman Bharat Digital Mission building foundational digital health infrastructure. IndiaA...
\item[Thailand] Asia; Developing; overall 3.9/10; Thai FDA (Food and Drug Administration); 15 devices; first AI instrument 2022. Scores: P5 C4 A4 T4 E4 PM3 L3. Privacy: PDPA (Personal Data Protection Act, 2022). AI rule: National AI Strategy and Action Plan (2022-2027). No binding AI law.. Challenges: Regulatory framework for AI devices still developing. Digital health infrastructure building. Balancing medical tourism innovation with domestic healthcare p... Notable: PDPA implementation is a major privacy milestone. Thailand 4.0 drives digital health investment. Medical tourism sector catalyzing AI adoption.
\item[South Africa] Africa; Emerging; overall 3.6/10; SAHPRA (South African Health Products Regulatory Authority); 10 devices; first AI instrument 2021. Scores: P6 C3 A3 T3 E4 PM3 L3. Privacy: POPIA (Protection of Personal Information Act, 2021). AI rule: No AI-specific regulation. AI Policy Framework under development.. Challenges: SAHPRA capacity and resource constraints. Digital infrastructure gaps. Brain drain in health AI expertise. Balancing innovation with addressing basic healthc... Notable: POPIA is Africa's most comprehensive data protection law. CSIR and universities driving AI health research. Growing health tech ecosystem in Cape Town and Jo...
\item[Kenya] Africa; Early; overall 2.7/10; PPB (Pharmacy and Poisons Board); 3 devices; first AI instrument 2019. Scores: P5 C2 A3 T2 E3 PM2 L2. Privacy: Data Protection Act (2019). AI rule: No AI-specific regulation. Blockchain and AI Taskforce Report (2019).. Challenges: Limited regulatory capacity. Infrastructure constraints. Prioritizing basic healthcare access alongside innovation. Dependency on international standards and... Notable: Kenya is a mobile health innovation leader in Africa. Data Protection Act among earliest in East Africa. Growing AI-for-health research at universities and t...
\item[Nigeria] Africa; Early; overall 2.6/10; NAFDAC (National Agency for Food and Drug Administration and Control); 5 devices; first AI instrument 2023. Scores: P4 C2 A3 T2 E3 PM2 L2. Privacy: NDPR (Nigeria Data Protection Regulation, 2019) / NDPA (Nigeria Data Protection Act, 2023). AI rule: National AI Strategy (draft). No binding AI regulation.. Challenges: Limited regulatory capacity for AI devices. Digital infrastructure gaps. Healthcare access challenges precede AI adoption. Funding constraints for regulatory... Notable: NDPA 2023 is a significant step in data governance. Growing AI research community (e.g., Data Science Nigeria). Mobile health innovations bypassing tradition...
\end{description}
}